% Gerado por etapa9/publicar_tabelas.py; fontes e células registradas na auditoria.
\begin{table}[H]
\centering
\normalsize
\setlength{\tabcolsep}{2pt}
\renewcommand{\arraystretch}{1}
\normalcaption\captionstyle{\centering}
\caption{Contrastes de capacidade e perda para aumento de 10\% no preço (p.p.)}
\label{tab:h1-magnitudes}
\begin{tabular}{p{2.4cm}*{4}{p{3.1cm}}}
\hline
Cultura / hipótese & Estimativa & IC95 mun. & IC95 mun.+ano & IC95 HAC(2) \\
Soja / H1a & 0,00385 & [-0,00618; 0,01388] & [-0,00666; 0,01437] & [-0,00273; 0,01044] \\
Soja / H1b & 0,00404 & [-0,00157; 0,00964] & [-0,00452; 0,01259] & [-0,00517; 0,01324] \\
Milho / H1a & 0,00307 & [-0,00036; 0,00650] & [-0,00139; 0,00752] & [-0,00075; 0,00688] \\
Milho / H1b & -0,00013 & [-0,00231; 0,00206] & [-0,00211; 0,00186] & [-0,00110; 0,00085] \\
\hline\end{tabular}
\par\smallskip
\begin{minipage}{\textwidth}\normalsize\centering
Fonte: elaboração própria com MapBiomas 10.1, IBGE/PAM, GAEZ v4, Banco Mundial, CPI-U e CHIRPS.
\par\medskip
\raggedright Legenda: H1a: um desvio-padrão adicional de capacidade; H1b: dez pontos percentuais adicionais de perda. p.p.: pontos percentuais. IC95: intervalo individual de 95\%. mun.: agrupamento municipal; mun.+ano: agrupamento duplo; HAC(2): covariância robusta à heterocedasticidade e à autocorrelação, banda dois.
\end{minipage}
\end{table}
